\clearpage
\sscp{\tsc{Flores-Lembata}}{03-02-04}
The \tsc{Flores-Lembata} node contains Sika,
spoken in eastern Flores,
the \tsc{Lamaholot} languages, spoken in far-eastern Flores
and the Solor archipelago (with coastal pockets on Pantar
and Alor represented by Alorese [aol]),
and Kedang spoken on the eastern-most peninsula of Lembata Island.
The history of the \tsc{Flores-Lembata} languages has been investigated
in detail by \citet{Fricke2019}
and we draw extensively on that work for our discussion.
\citet{Fricke2019} also contains a grammar sketch of Kalikasa (a.k.a.
Central Lembata; ISO [lvu]) which is the first detailed
description a \tsc{Central Lamaholot} language.\footnote{ 
		\citet{Krauße2016} also provides a short description of \tsc{Central
		Lamaholot} Painara (a.k.a.
		South Lembata, Atadei ISO [lmf]).
		\citet{Fricke2015b} provides a word list of Painara.}
This node is conservative in several respects compared with other
\tsc{Bima-Lembata} languages.

Phonological evidence for \tsc{Flores-Lembata} comes from shared
*s > **h in a subset of words,
*ŋ > \it{n} word initially,
and merger of *d/*j/*z.
Although these sound changes are attested to some extent in other
branches of \tsc{Bima-Lembata}, no other branch is known to have all three sound changes.

The first sound change shared by \tsc{Flores-Lembata} is *s >
**h in a subset of words,
with subsequent **h > {\0} in some languages.
Examples are given \trf{03-41},
along with cognates in languages from each other branch of
\tsc{Bima-Lembata} to show that,
while this sound change has occurred in some languages -- notably
several languages of \tsc{Sumba-Havu} -- it cannot be posited
as a defining sound change for any other branch.

\begin{table}\tcp{\tsc{Flores-Lembata} with *s > **h}{03-41}
\begin{adjustbox}{width=\textwidth}
	\begin{threeparttable}\st{2pt}
	\begin{tabular}{lllllllll}\lsptoprule
*gloss	&	dog	&	bow	&	climb	&	one, alone	&	heart	&	salt	&	thin	&	spouse\\
PMP	&	*a\tbr{s}u	&	*bu\tbr{s}uR	&	*\tbr{s}akay	&	*\cg{(ma)}-ə\tbr{s}a	&	*pu\tbr{s}uŋ\su{Δ}	&	*qa\tbr{s}iRa	&	*ma-nipi\tbr{s}	&	*qa\tbr{s}awa\\	\midrule
Sika	&	{\hr}a\tbr{h}u	&	{\hr}βu\tbr{h}ur	&	{\hr}\tbr{h}aʔe	&	{\hr}me\tbr{h}a	&	{\hr}pu\tbr{h}uŋ	&	\cg{(hini)}	&	\cg{(lata)}	&	\\
Ile Mandiri	&	{\hr}a\tbr{h}o	&	{\hr}vu\tbr{h}uʔ	&	{\hr}\tbr{h}aka\su{‡}	&	{\hr}me\tbr{h}aʔ	&	{\hr}pu\tbr{h}un	&	{\hqa}\tbr{s}iʔa\su{ǁ}	&	{\hr}mənipi	&	\\
Lamalera	&	{\hr}ao	&		&	{\hr}\tbr{h}aka\su{‡}	&		&	{\hr}puo	&	{\hqa}\tbr{s}ia\su{ǁ}	&	{\hr}mənipi	&	\\
Alorese	&	{\hr}a\tbr{h}o\su{†}	&	{\hr}wu\tbr{h}u	&		&		&	\cg{(kubaŋ)}	&	{\hqa}\tbr{s}iʔa, \tbr{s}ia\su{ǁ}	&	{\hr}mənipi	&	\\
Kalikasa	&	{\hr}au	&	{\hr}vuvu	&	{\hh}akaʤ	&	{\hr}mea	&	{\hr}puo	&	\hp{*qas}ira	&	{\hr}mipiv	&	{\hh}ava-n\\
Painara	&	{\hr}ao-r	&	\hp{*v}uvo-r	&		&		&	{\hr}puo-n	&	\hp{*qas}ira-r	&	{\hr}mipivən	&	\\
Lamatuka	&	{\hr}a\tbr{h}o	&		&		&		&	{\hr}pu\tbr{h}o	&	{\hqa}\tbr{h}ira	&	{\hr}mipi\tbr{h}i	&	{\hr}\tbr{h}awã\\
Lewo Eleng	&	{\hr}a\tbr{h}o	&		&		&		&	{\hr}pũ\tbr{h}ũ	&	\hp{qa}\cg{(taʔo)}	&	{\hr}mipi	&	\\
Kedang	&	{\hr}au	&	{\hr}vur	&	{\hh}aʔ\su{‡}	&	{\hr}e\tbr{h}aʔ	&	{\hr}pu\tbr{h}un	&	\hp{qa}\cg{(tæʔu)}	&	{\hr}mipi	&	\\	\midrule
Bima	&	\cg{(lako)}	&	\cg{(ina fana)}	&	\cg{(teka)}	&	{\hr}iʧa	&	\cg{(hodo)}	&	{\hqa}sia	&	{\hr}nipi	&	\\
Manggarai	&	{\hr}aʧu	&	{\hr}ʋusu	&	\cg{(tuke)}	&	{\hr}ʧa	&	{\hr}puʧu	&	{\hqa}ʧiʔe	&	{\hr}mipis	&	\\
Mamboru	&	{\hr}asu	&	{\hr}lawosur-a	&	{\hr}sawi	&	{\hr}mesa	&	{\hr}kapusu	&	\hp{qa}\cg{(masi)}	&	{\hr}manipis-a	&	\\
		\lspbottomrule
	\end{tabular}
		\begin{tablenotes}\tnsize
			\item[†]	{Both \it{aho}
								and \it{asu} ‘dog’ are given for Alorese in Bana village by \citet{Sulistyono2018}.
								The second word is presumably a borrowing given it has both irregular
								*s = \it{s} and *u = \it{u} /{\gap}{\#}.}
			\item[‡]	{Ile Mandiri \it{haka} is ‘come from a lower level,
								come westwards/from the sea’,
								Lamalera \it{haka} ‘east’,
								Kedang \it{aʔ} is ‘from,
								hither, come from, an approximately level direction seawards
								(or westwards) towards the speaker’.}
			\item[Δ]	{The Ile Mandiri,
								Kedang and Sika reflexes mean ‘flower, blossom’.
								Sika also has \it{puhuŋ wuaŋ} ‘heart’.}
			\item[{ǁ}]	{\tsc{West Lamaholot}
								and Alorese \it{siʔa, sia} ‘salt’ are possibly secondary borrowings.
								Salt was a historic trade commodity in this area.
								Non-\tsc{An} \tsc{Alor-Pantar} languages
								also have borrowings of *qasiRa ‘salt’.}
		\end{tablenotes}
	\end{threeparttable}
\end{adjustbox}
\end{table}

We cannot posit *s > **h in all cases at the level of Proto-Flores-Lembata
as several \tsc{Central Lamaholot} languages -- Kalikasa,
Levuka, Lewokukun, and Mingar -- have retained *s = \it{s} in some words.
Similarly, Ile Mandiri Lamaholot (Lewolema dialect) is known
to have *s > \it{h {\tl} s} in some of these same words.
Where cognates are known in other members of \tsc{Flores-Lembata}
they have *s > \it{h} (> {\0}),
though more detailed data may reveal additional cases of *s > \it{h {\tl} s}.
Examples are given in \trf{03-42},
with retentions of *s = \it{s} shaded.

\begin{table}\tcp{\tsc{Flores-Lembata} with *s = **s}{03-42}
\begin{adjustbox}{width=\textwidth}
	\begin{threeparttable}\st{2pt}
	\begin{tabular}{llllllll}\lsptoprule
local gloss	&	wash	&	mouth	&	sew	&	breast	&	navel	&	nine	&	wrong\\
PMP	&	*ba\tbr{s}əq	&	*ŋu\tbr{s}uq	&	**\tbr{s}auR\su{†}	&	*su\tbr{s}u\su{‡}	&	*pu\tbr{s}əj	&	*\tbr{s}iwa	&	*\tbr{s}alaq\\ \midrule
Sika	&	\cg{(popo)}	&	\cg{(βaaŋ)}	&	{\hr}\tbr{h}aut	&	{\hr}u\tbr{h}uŋ	&	{\hr}pu\tbr{h}ər	&	{\hr}\tbr{h}iβa	&	{\hr}\tbr{h}ala\\
Ile Mandiri	&	{\hr}ba\tbr{h}aʔ	&	{\hr}nu\tbr{h}u, nu\tbr{s}u	{\cgr}&	{\hr}\tbr{h}avu	&	{\hr}tu\tbr{h}o, tu\tbr{s}o	{\cgr}&	{\hr}kəpu\tbr{h}u\cg{(r)}	&	{\hr}\tbr{h}iva	&	{\hr}\tbr{h}alaʔ\su{Δ}\\
Lamalera	&	{\hr}baa	&	\cg{(fəfã)}	&	{\hr}\tbr{h}au	&	{\hr}tuo	&	{\hr}kəpuur	&	{\hr}\tbr{h}ifa	&	\cg{(əla)}\\
Alorese	&	\cg{(bema)}	&	{\hr}nu\tbr{h}uŋ	&	{\hr}\tbr{h}aur	&	{\hr}tuho	&	{\hr}pu\tbr{h}or	&	{\hr}\tbr{h}iwa	&	{\hr}\tbr{h}ala\\
Kalikasa	&	{\hr}ba\tbr{s}a	{\cgr}&	{\hr}nu\tbr{s}ən	{\cgr}&	{\hr}\tbr{s}aur	{\cgr}&	{\hr}tu\tbr{s}u	{\cgr}&	{\hr}kəpu\tbr{s}ər	{\cgr}&	{\hr}\tbr{s}iva	{\cgr}&	{\hr}\tbr{s}\<n\>alak{\cgr}\\
Levuka	&	{\hr}ba\tbr{s}a	{\cgr}&	{\hr}nu\tbr{s}	{\cgr}&	{\hr}\tbr{s}aur	{\cgr}&		&		&		&	\\
Mingar	&	\cg{(pusɤ)}	&	{\hr}nu\tbr{s}a	{\cgr}&	{\hr}\tbr{s}aur	{\cgr}&		&		&		&	\\
Lewokukun	&	{\hr}ba\tbr{s}a	{\cgr}&	{\hr}nu\tbr{s}ə-ga	{\cgr}&	{\hr}\tbr{s}aur	{\cgr}&		&		&		&	\\
Labalekan	&	{\hr}ba\tbr{h}a	&	\cg{(fəfaha)}	&	{\hr}\tbr{h}aur	&		&		&		&	\\
Painara	&	{\hr}ba\tbr{h}aŋ	&	{\hr}nu\tbr{h}ə	&	{\hr}\tbr{h}aurəŋ	&	{\hr}tu\tbr{h}o-r	&	{\hr}kəpu\tbr{h}ər	&	{\hr}\tbr{h}iva	&	\cg{(bənələŋən)}\\
Lamatuka	&	{\hr}ba\tbr{h}a	&	{\hr}nu\tbr{h}e	&	{\hr}\tbr{h}eu	&		&		&		&	\\
Lewo Eleng	&	{\hr}ba\tbr{h}a	&	{\hr}nu\tbr{h}e	&	{\hr}\tbr{h}eu	&		&		&		&	\\
Kedang	&	{\hr}ba\tbr{h}iŋ	&	\cg{(nunu)}	&	{\hr}\tbr{h}æwuŋ	&	{\hr}tu	&	{\hr}pu\tbr{h}e	&	\cg{(leme apaʔ)}	&	{\hr}alen\\
		\lspbottomrule
	\end{tabular}
		\begin{tablenotes}\tnsize
			\item[†]	{For **sauR compare Proto-Rote-Meto *soo ‘sew’ \citep[340]{Edwards2021},
								Kisar \it{hour} ‘sew’, Uruangnirin \it{a-sor} ‘sew’,
								Proto-Central-Maluku **sa(q)u \citep{Collins1983}, and PCMP **sora \citep{BlustTrussel2020}.}
			\item[‡]	{PMP *susu ‘breast’ has irregular initial *s > \it{t} in \tsc{Flores-Lembata}.}
			\item[Δ]	{Ile Mandiri \it{halaʔ} is ‘no, not’.
								It has \it{nalan} < **s\<n\>ala for ‘wrong’.}
		\end{tablenotes}
	\end{threeparttable}
\end{adjustbox}
\end{table}

The data relating to PMP *s can be explained by positing that
*s > **h was a sound change in progress which had not completely
diffused through the lexicon by the time of the break-up of Proto-Flores-Lembata.
After the break-up of the proto-language,
this sound change continued in most languages except some members
of \tsc{Central Lamaholot} and \tsc{West Lamaholot}.

Another sound change shared by \tsc{Flores-Lembata} languages is word
initial *ŋ > \it{n} in at least four words.
Examples are given in \trf{03-43}.
Of the words given in \trf{03-43},
we can probably posit universal *ŋ > \it{n} for the first four
words (though the reflexes of *ŋilu are ambiguous in Alorese and Kedang).
However, *ŋ = \it{ŋ} word medially in *ŋaŋa ‘open mouth’,
as well as *sa-ŋa-Ratus > Sika \it{ŋasu} ‘hundred’.
An additional word which retains initial *ŋ = \it{ŋ} in Sika
is *ŋəŋəŋ > \it{ŋeŋuŋ} ‘hum’.

Of the words with *ŋ > \it{n} in \tsc{Flores-Lembata},
only *ŋajan and *ŋuda set these languages apart from other \tsc{Bima-Lembata}
languages, as other \tsc{Bima-Lembata} languages retain *ŋ = \it{ŋ} in these words.
Examples include: *ŋajan ‘name’ > Bima,
Ende, Kambera, and Dhao \it{ŋara} ‘name’,
Manggarai \it{ŋasaŋ} ‘name’, as well as *ŋuda ‘young (plants)’
> Ende \it{ŋura}, Kambera \it{kaŋura}, Dhao \it{ŋəru}.
Reflexes of *ŋusuq ‘mouth’ are not known in other \tsc{Bima-Lembata}
languages, while *ŋ > \it{n}
in some languages of Flores for reflexes of
*ŋilu ‘sour’, such as Ende \it{ni\,ɹu}.\footnote{
		Other known reflexes of *ŋilu ‘sour’ include Kambera,
		Lewa, Mangili \it{yilu}; Anakalang, Mamboru \it{mayilu};
		and Havu \it{meɲilu}.}

\largerpage
\begin{table}\tcp{\tsc{Flores-Lembata} lexically specific *ŋ > \it{n} /{\#\gap}}{03-43}
\begin{adjustbox}{width=\textwidth}
	\begin{threeparttable}\st{2pt}
	\begin{tabular}{lllllll}\lsptoprule
*gloss	&	name	&	young (plants)	&	snout	&	sour	&	open mouth	&	hundred\\
PMP	&	*\tbr{ŋ}ajan	&	*\tbr{ŋ}uda	&	*\tbr{ŋ}usuq	&	*\tbr{ŋ}ilu\su{†}	&	*\tbr{ŋ}aŋa	&	*sa-\tbr{ŋ}a-Ratus\\ \midrule
Sika	&	{\hr}\tbr{n}aran	&	{\hr}\tbr{n}urak	&	\cg{(βaaŋ)}	&	{\hr}\tbr{n}iluk	&	{\hr}\tbr{ŋ}aŋa	&	{\hr}\tbr{ŋ}asu\\
Ile Mandiri	&	{\hr}\tbr{n}aran	&	{\hr}mə\tbr{n}urenʔ	&	{\hr}\tbr{n}uhu, \tbr{n}usu	&	{\hr}kə\tbr{n}ilun	&	{\hr}\tbr{n}aŋa	&	{\hr}ratu\\
Lamalera	&	{\hr}\tbr{n}araŋ	&	{\hr}mə\tbr{n}ura-	&	\cg{(fəfã)}	&		&		&	{\hr}ratu\\
Alorese	&	{\hr}\tbr{n}araŋ	&		&	{\hr}\tbr{n}uhuŋ	&	{\hr}gilo	&		&	{\hr}ratu\\
Kalikasa	&	{\hr}\tbr{n}aʤan	&		&	{\hr}\tbr{n}usən	&	{\hr}k\tbr{n}iluk	&		&	{\hr}ratu tu\\
Painara	&	{\hr}\tbr{n}ayan	&		&	{\hr}\tbr{n}uhə	&		&		&	{\hr}rat\\
Lamatuka	&	{\hr}\tbr{n}arã	&		&	{\hr}\tbr{n}uhe	&		&		&	\\
Lewo Eleng	&	{\hr}\tbr{n}ara	&		&	{\hr}\tbr{n}uhe	&		&		&	\\
Kedang	&	{\hr}\tbr{n}aya	&	{\hr}\tbr{n}uyan	&	\cg{(nunu)}	&	{\hr}kiru	&	{\hr}maŋe	&	{\hr}ratuʔ\\
		\lspbottomrule
	\end{tabular}
		\begin{tablenotes}\tnsize
			\item[†]	{
								Reflexes of *ŋilu appear to be from **ka-ŋilu > **kŋilu and/or
								**knilu with subsequent simplification of the consonant cluster.
								Alorese and Kedang appear to retain initial **k with deletion
								of the nasal.}
		\end{tablenotes}
	\end{threeparttable}
\end{adjustbox}
\end{table}

Another change probably shared by members of \tsc{Flores-Lembata}
is merger of *d/*j/*z in most cases.
The reflexes are somewhat complex
and require detailed discussion.
Word medially, the reflexes are mostly clear: \tsc{Central Lamaholot}
has *j/*z/*d > **ʤ (with sporadic **ʤ > \it{y} in Painara).
Kedang usually has *j/*z/*d > \it{y.} Other \tsc{Flores-Lembata}
languages have *j/*z/*d > \it{r}.
Examples are given in \trf{03-44}.
An additional example of medial *d can be seen in reflexes of
*ŋuda ‘young (plants)’ in \trf{03-43}.

Although Kedang usually has *j/*z/*d > \it{y},\footnote{\label{fn:Kedangd/z}
		Three additional examples of Kedang *d/*z > \it{y} /V{\gap}V
		are *kudən > \it{uyeŋ} ‘pot’,
		ma-tazəm > \it{dæyæʔ} ‘sharp’ (Sika \it{diran} ‘sharp’ may also
		be from *ma-tazəm with irregular vowels),
		and *zəlay > \it{leye} ‘Job’s tears’ with metathesis of the initial
		and medial consonant before *z > \it{y}.}
it also has
some cases of *j/*d > \it{r}.
Four examples have been identified: *qudaŋ > \it{uraŋ} ‘shrimp’,
*udu > \it{uru} ‘grass’,
*huaji > \it{ariʔ} ‘younger sibling’,
and *bujəq > \it{vuran} ‘foam’.
We have not yet identified any examples of intervocalic *z > \it{r} in Kedang.
Similarly, the reflexes of *pija have irregular *j > \it{r} in
\tsc{Central Lamaholot} languages, exemplified by Kalikasa and Painara \it{pira} ‘how much?’.\footnote{Kedang
		has *pija › \it{pie} with presumable *j > **y > {\0} /i{\gap}
		--- compare *ŋijuŋ > \it{niŋ} ‘nose’ in \trf{03-44}.
		\tsc{West Lamaholot} and Alorese have *pija > \it{pira} ‘how
		much?’ with regular *j > \it{r}.}

\begin{table}\tcp{\tsc{Flores-Lembata} *d/*z/*j /V{\gap}V}{03-44}
\begin{adjustbox}{width=\textwidth}
	\begin{threeparttable}\st{2pt}
	\begin{tabular}{llllllll}\lsptoprule
local gloss	&	rain	&	spit	&	sun	&	name	&	nose	&	alive	&	white\\
PMP	&	*qu\tbr{z}an	&	*qi\tbr{z}uR\su{†}	&	*qalə\tbr{j}aw	&	*ŋa\tbr{j}an	&	*\cg{(ŋ)}i\tbr{j}uŋ	&	*ma-qu\tbr{d}ip	&	*bu\tbr{d}aq\\	\midrule
Sika	&	{\hq}u\tbr{r}aŋ	&	{\hr}ni\tbr{r}u	&	{\hqa}lə\tbr{r}o	&	{\hr}na\tbr{r}an	&	{\hr}i\tbr{r}u	&	{\hr}mo\tbr{r}et	&	{\hr}bu\tbr{r}a\\
Ile Mandiri	&	{\hq}u\tbr{r}an	&	{\hr}pə\tbr{r}ino \it{(met.)}	&	{\hqa}rə\tbr{r}a	&	{\hr}na\tbr{r}an	&	{\hr}i\tbr{r}u	&	{\hr}mo\tbr{r}it	&	{\hr}bu\tbr{r}aʔ\\
Adonara	&	{\hq}u\tbr{r}ã	&	{\hr}pə\tbr{r}ino \it{(met.)}	&	{\hqa}lə\tbr{r}a, rə\tbr{r}a	&	{\hr}na\tbr{r}ã	&	{\hr}i\tbr{r}ũ	&	{\hr}mo\tbr{r}ip, mo\tbr{r}it	&	{\hr}bu\tbr{r}aʔ\\
Lamalera	&	{\hq}u\tbr{r}ã	&	{\hr}təmi\tbr{r}o	&	{\hqa}lə\tbr{r}a	&	{\hr}na\tbr{r}aŋ	&	{\hr}ni\tbr{r}uŋ	&	{\hr}mo\tbr{r}i	&	{\hr}bu\tbr{r}ã\\
Alorese	&	{\hq}u\tbr{r}aŋ	&	{\hr}pəni\tbr{r}o, ni\tbr{r}o	&	{\hqa}lə\tbr{r}a	&	{\hr}na\tbr{r}aŋ	&	{\hr}i\tbr{r}uŋ, ni\tbr{r}uŋ	&	{\hr}mo\tbr{r}ik	&	{\hr}bu\tbr{r}ak\\
Kalikasa	&	{\hq}u\tbr{ʤ}an	&	{\hr}təmi\tbr{ʤ}u	&	\hp{qa}\cg{(luvak)}	&	{\hr}na\tbr{ʤ}an	&	{\hr}ni\tbr{ʤ}u	&	{\hr}mo\tbr{ʤ}ip	&	{\hr}bu\tbr{ʤ}ak\\
Painara	&	{\hq}u\tbr{ʤ}an	&	{\hr}təmi\tbr{ʤ}o	&	{\hqa}lə\tbr{ʤ}aw	&	{\hr}na\tbr{y}an	&	{\hr}ni\tbr{ʤ}u	&	{\hr}mo\tbr{ʤ}ipən	&	{\hr}bu\tbr{ʤ}ak\\
Lamatuka	&	{\hq}u\tbr{r}ã	&	{\hr}pi\tbr{r}u	&	\hp{qa}\cg{(luwa)}	&	{\hr}na\tbr{r}ã	&	{\hr}i\tbr{r}u	&	{\hr}mo\tbr{r}i	&	{\hr}bu\tbr{r}ʔã\\
Lewo Eleng	&	{\hq}u\tbr{r}ã	&	{\hr}pi\tbr{r}u	&	{\hqa}lə\tbr{r}a	&	{\hr}na\tbr{r}a	&	{\hr}ni\tbr{r}ũ	&	{\hr}mo\tbr{r}i	&	{\hr}bu\tbr{r}õ\\
Kedang	&	{\hq}u\tbr{y}a	&	{\hr}i\tbr{y}uʔ, mi\tbr{y}uʔ	&	{\hqa}lo\tbr{y}o	&	{\hr}na\tbr{y}a	&	{\hr}niŋ\su{‡}	&	\cg{(bita)}	&	{\hr}bu\tbr{y}aʔ\\
		\lspbottomrule
	\end{tabular}
		\begin{tablenotes}\tnsize
			\item[†]	{Reflexes of *qizuR ‘spit’ usually have initial elements which
								may be fossilised prefixes.}
								%Initial \it{tə-} may be a reflex of *taR- ‘spontaneous action’.
								%For other words in the region with initial \it{m} compare Proto-Rote-Meto
								%*midu and Onin \it{miri{\tl}miri}.
								%The word \it{ilu} ‘spittle’ occurs in a number of \tsc{Bima-Lembata}
								%languages including Ile Mandiri,
								%Lamalera, Alorese, and Havu.
								%This word may also be from *qizuR.}
			\item[‡]	{Kedang \it{niŋ} ‘nose’ may have *j > **y > {\0} /i{\gap},
								but *qizuR > \it{iyuʔ} ‘spit’ shows this is not a regular change.}
		\end{tablenotes}
	\end{threeparttable}
\end{adjustbox}
\end{table}

Word initially, *z and *d also merge in most cases in \tsc{Flores-Lembata}
(initial *j does not occur).
Examples are given in \trf{03-45}.
The clearest cases are the first five cognate sets which have
*z/*d > \it{r} in Sika,
and *z/*d > \it{d} in all other languages,
though *dəpa ‘arm span’ displays a split in Ile Mandiri between
*d = \it{d} in the verbal form
and *d > \it{r} in the nominal form.\footnote{
		A similar split between reflexes of *dəpa is also seen in the
		languages of Sumba, as exemplified by Kambera \it{rɐpa} ‘arm span’
		%with *d > \it{r}
		but \it{ndɐpa} ‘arm span,
		measure by arm spans’.}
%*d > \it{nd}.}
We can thus posit
initial *z/*d > **d at the level of Proto-Flores-Lembata
for these five words (excepting the verbal reflexes of *dəpa),
with subsequent **d > \it{r} in Sika.

\begin{table}\tcp{\tsc{Flores-Lembata} *d/*z /{\#\gap}}{03-45}
\begin{adjustbox}{width=\textwidth}
	\begin{threeparttable}\st{2pt}
	\begin{tabular}{llllll|lll}\lsptoprule
local gloss	&	hear	&	stand	&	far	&	bad	&	arm span	&	path	&	bone	&	two\\
PMP	&	*\tbr{d}əŋəR	&	*\tbr{d}iRi	&	*\tbr{z}auq	&	*\tbr{z}aqat	&	*\tbr{d}əpa	&	*\tbr{z}alan	&	*\tbr{d}uRi	&	*\tbr{d}uha\\	\midrule
Sika	&	{\hr}\tbr{r}əna	&	\cg{(gəra)}	&	\cg{(blaβir)}	&	\cg{(goʔis)}	&	{\hr}\tbr{r}əpa	&	{\hr}\tbr{l}alaŋ	&	{\hr}\tbr{l}uriŋ	&	{\hr}\tbr{r}ua\\
Sika, Hewa	&	{\hr}\tbr{r}əna	&	\cg{(ʔ-əra)}	&	\cg{(blavir)}	&	\cg{(goʔit)}	&		&	{\hr}la\tbr{r}an	&	{\hr}\tbr{l}uri	&	{\hr}\tbr{r}ua\\
Ile Mandiri	&	{\hr}\tbr{d}əŋəʔ	&	{\hr}\tbr{d}eʔi, \tbr{d}eʔin	&	{\hr}\tbr{d}oan	&	{\hr}\tbr{d}a, \tbr{d}ate	&	{\hr}\tbr{r}əpa,\tbr{d}əpan\su{†}	&	{\hr}ra\tbr{r}an	&	{\hr}\tbr{r}iʔuk	&	{\hr}\tbr{r}ua\\
Adonara	&	{\hr}\tbr{d}əŋəʔ	&	{\hr}\tbr{d}eʔi	&	{\hr}\tbr{d}oã	&	\cg{(mədõ)}	&		&	{\hr}la\tbr{r}ã,ra\tbr{r}ã	&	{\hr}\tbr{r}iʔũ	&	{\hr}\tbr{r}ua\\
Lamalera	&	{\hr}\tbr{d}əŋa	&	{\hr}\tbr{d}ei	&	{\hr}\tbr{d}oe	&	{\hr}afa\tbr{d}atəŋ	&		&	{\hr}la\tbr{r}ã	&	{\hr}\tbr{r}iuk	&	{\hr}\tbr{r}ua\\
Alorese	&	{\hr}\tbr{d}əŋa	&	\cg{(tide)}	&	{\hr}\tbr{d}oaŋ	&	{\hr}\tbr{d}ate	&		&	{\hr}la\tbr{r}aŋ	&	{\hr}\tbr{r}uiŋ	&	{\hr}\tbr{r}ua\\
Kalikasa	&	{\hr}\tbr{d}əŋər	&	{\hr}\tbr{d}iri	&	{\hr}\tbr{d}oak	&	{\hr}kə\tbr{d}atək	&		&	{\hr}\tbr{l}alan	&	{\hr}\tbr{r}iuk	&	{\hr}\tbr{ʤ}ua\\
Painara	&	{\hr}\tbr{d}əŋə	&	{\hr}\tbr{d}ir	&	{\hr}\tbr{d}oakən	&		&		&	{\hr}\tbr{l}alan	&	{\hr}\tbr{r}iuk	&	{\hr}\tbr{ʤ}ua\\
Lamatu.	&	{\hr}\tbr{d}əŋe	&	{\hr}\tbr{d}iri	&	{\hr}\tbr{d}oa	&	{\hr}\tbr{d}atʔẽ	&		&	{\hr}la\tbr{r}ã	&	{\hr}\tbr{r}iʔũ	&	\\
Lewo El.	&	{\hr}\tbr{d}əŋe	&	{\hr}\tbr{d}iri	&	{\hr}\tbr{d}oa	&	{\hr}\tbr{d}atẽ	&		&	{\hr}la\tbr{r}an	&	{\hr}\tbr{r}iʔũ	&	\\
Kedang	&	{\hr}\tbr{d}æŋær	&	{\hr}ma\tbr{d}er	&	{\hr}\tbr{d}oa	&	{\hr}\tbr{d}aten, &	{\hr}\tbr{d}æpæ	&	{\hr}\tbr{l}ala	&	{\hr}\tbr{l}urin	&	{\hr}\tbr{s}ue\\
&&&&{\hr}\tbr{d}atoʔ	&&&&\\
		\lspbottomrule
	\end{tabular}
		\begin{tablenotes}\tnsize
			\item[†]	{Ile Mandiri \it{dəpan} is given as a verb ‘measure by arm spans’ (Indonesian \it{mendepa}).}
		\end{tablenotes}
	\end{threeparttable}
\end{adjustbox}
\end{table}

The final three words in \trf{03-45} show other
developments for initial *d/*z.
The reflexes of *zalan in Flores-Lembata probably have metathesis
of the initial and medial consonants: *zalan > **lazan with subsequent assimilation
of the medial consonant to \it{l} /lV{\gap} in Sika,
\tsc{Central Lamaholot}, and Kedang.
Assimilation has not taken place in \tsc{East Lamaholot}
and \tsc{West Lamaholot}.\footnote{ 
		Ile Mandiri (Lewolema dialect) \it{raran} ‘path’ has regressive
		assimilation of initial **l > \it{r} /{\gap}Vr.
		This is also seen in *qaləjaw > **ləra › \it{rəra} (see \trf{03-44}).}

The reflexes of *duRi have several complications.
Apart from Sika and Kedang, all putative reflexes require positing
vowel metathesis and irregular *R > **ʔ,
{\0} in \tsc{Central Lamaholot} and \tsc{East Lamaholot}.
(*R > **ʔ, {\0} is regular in \tsc{West Lamaholot}.) It may be
the case that these words are not in fact reflexes of *duRi.
Sika and Kedang have irregular initial *d > \it{l},
which may be due to sporadic dissimilation from the following
liquid *R = \it{r}.
The reflexes of *duha would be regular for medial *d,
except for Kedang which has irregular *d > \it{s}.

To summarise, *z, *j, and *d merge in most words in \tsc{Flores-Lembata}.
\citet{Fricke2019} posits that the outcome was Proto-Flores-Lembata
**d in all environments with this reflex retained word initially
in all branches except Sika,
which has **d > \it{r}.
Under this hypothesis Sika,
\tsc{West Lamaholot}, and \tsc{East Lamaholot} have medial **d
> \it{r}, while \tsc{Central Lamaholot} has **d > \it{ʤ}
and Kedang has **d > \it{r,
y} -- with \it{y} probably developing from intermediate **ʤ.
\tsc{Central Lamaholot} and Kedang thus potentially share **d > **ʤ.

\largerpage
However, an alternate interpretation of the data would be to
posit medial *d/*z/*j > Proto-Flores-Lembata **ʤ,
with subsequent **ʤ > \it{r} in \tsc{West Lamaholot} and \tsc{East Lamaholot}.
If this analysis is adopted,
the reflexes in \tsc{Central Lamaholot}
and Kedang would be retentions
and would not provide any subgrouping evidence.

Within \tsc{Flores-Lembata} there is good phonological support
for a \tsc{West Lamaholot} node based on medial *R > **ʔ,
with subsequent **ʔ > {\0} in some varieties.
Other \tsc{Flores-Lembata} languages usually have *R = \it{r}.
If PMP *R is taken to be alveolar [r],
and the outcome of *d/*z/*j word initially were **d,
then the changes affecting this proto-phoneme in \tsc{West Lamaholot}
must have occurred before the shift of *z/*j/*d > \it{r} in \tsc{West Lamaholot}.
Examples of medial *R are given in \trf{03-46},
where shading indicates shared *R > **ʔ (> {\0}).

\begin{table}\tcp{\tsc{Flores-Lembata} *R}{03-46}
\begin{adjustbox}{width=\textwidth}
	\st{2pt}\begin{tabular}{llllllll}\lsptoprule
local gloss	&	new	&	heavy	&	stand	&	red	&	flee	&	turtle	&	dry in sun\\
PMP	&	*baqə\tbr{R}u	&	*bə\tbr{R}əqat	&	*di\tbr{R}i	&	*ma-i\tbr{R}aq	&	*pa-la\tbr{R}iw	&	**ke\tbr{R}a	&	*pa-wa\tbr{R}i\\	\midrule
Sika	&	{\hr}βə\tbr{r}u	&	{\hr}bə\tbr{r}at	&	\cg{(gəra)}	&	{\hr}me\tbr{r}ak	&	{\hr}pla\tbr{r}i	&	\hp{**}ʔe\tbr{r}a	&	{\hr}ʔβo\tbr{r}i\\
\rcl	Ile Mandiri	&	{\hr}vu\tbr{ʔ}u	&	{\hr}ba\tbr{ʔ}at	&	{\hr}de\tbr{ʔ}i	&	{\hr}me\tbr{ʔ}a	&	{\hr}pəla\tbr{ʔ}e	&	\hp{**}ke\tbr{ʔ}a	&	{\hr}pa\tbr{ʔ}in\\
\rcl	Adonara	&	{\hr}wu\tbr{ʔ}ũ	&	{\hr}ba\tbr{ʔ}at	&	{\hr}de\tbr{ʔ}i	&	{\hr}me\tbr{ʔ}a	&	{\hr}pala\tbr{ʔ}e	&	\hp{**}ke\tbr{ʔ}a	&	{\hr}pa\tbr{ʔ}ĩ\\
\rcl	Lebatukan 	&	{\hr}wu\tbr{ʔ}ũ	&	{\hr}ba\tbr{ʔ}at	&	{\hr}de\tbr{ʔ}i	&	{\hr}me\tbr{ʔ}ã	&		&		&	\\
\rcl	Lamalera	&	{\hr}fuu	&	{\hr}batã	&	{\hr}dei	&	{\hr}meã	&	{\hr}pəlae	&	\hp{**}kea	&	\\
\rcl	Alorese	&	{\hr}wunoŋ	&	{\hr}baa	&	\cg{(tide)}	&	{\hr}meaŋ	&	{\hr}plae	&	\hp{**}kea	&	{\hr}paiŋ\\
Kalikasa	&	{\hr}və\tbr{r}un	&	{\hr}b‹ən›ə\tbr{r}at	&	{\hr}di\tbr{r}i	&	{\hr}me\tbr{r}an	&	{\hr}ka\tbr{r}i	&	\hp{**}kea-r	{\cgr}&	{\hr}pa\tbr{r}i\\
Painara	&	{\hr}və\tbr{r}un	&	{\hr}b‹ən›ə\tbr{r}ət-ən	&	{\hr}di\tbr{r}	&	{\hr}me\tbr{r}an	&	{\hr}ka\tbr{r}	&	\hp{**}ke\tbr{r}a-r	&	{\hr}pa\tbr{r}iŋ\\
Lamatuka	&	{\hr}wə\tbr{r}ũ	&	{\hr}bə\tbr{r}a	&	{\hr}di\tbr{r}i	&	{\hr}meã	{\cgr}&		&		&	\\
Lewo Eleng	&	{\hr}wə\tbr{r}ũ	&	{\hr}bə\tbr{r}a	&	{\hr}di\tbr{r}i	&	{\hr}me\tbr{r}a	&		&		&	\\
Kedang	&	{\hr}wæ\tbr{r}u	&	{\hr}ba\tbr{r}aʔ	&	{\hr}made\tbr{r}	&	{\hr}me\tbr{r}an	&	\cg{(bæyæŋ)}	&	\hp{**}ʔe\tbr{r}e	&	{\hr}pa\tbr{r}i daraŋ\\
		\lspbottomrule
	\end{tabular}
\end{adjustbox}
\end{table}

Word initially, there appears to be a split of *R > \it{r,} **ʔ/{\0}
in \tsc{West Lamaholot}, though only three examples of initial *R with reflexes have been identified.
*sa-ŋa-Ratus ‘hundred’ has *R > \it{r} in all known reflexes (see \trf{03-43}).
Alorese *Rumaq > \it{uma} has *R > **ʔ > {\0} (other \tsc{Flores-Lembata}
languages have no known reflex of *Rumaq).
PMP *Ramut ‘root’ has *R > *R > {\0}, \it{r}.
Examples include: Ile Mandiri,
Adonara \it{ramut {\tl} amut},
Lebatukan \it{ramu}, Lamalera \it{ramut}, and Alorese \it{ramuk}.
Sika has \it{ʔramut} which points to earlier **kramut.
This initial consonant cluster probably is the source of the
variable reflexes in \tsc{West Lamaholot}.\footnote{
		Initial consonant clusters also have variable reflexes in Ngadha.
		See the discussion in \citet[83--84]{Elias2018} for a succinct overview.}

There is evidence for placing \tsc{West Lamaholot},
\tsc{East Lamaholot}, and (perhaps) \tsc{Central Lamaholot} together
on the basis of lowering of word-final historical high vowels to mid.
Such vowel lowering only occurs in \tsc{Central Lamaholot} as
part of a word shape alternation in combination with a suffix
\it{-r} (see \citet[63ff]{Fricke2019} for discussion of this process
in Kalikasa), and does not occur in absolute word-final position.
Conversely, such vowel lowering is blocked in West Lamaholot
in some contexts when the vowel is not in absolute final position;
as shown by *hapuy ‘fire’ > **api > Ile Mandiri \it{ape} ‘fire’ but \it{api-n} ‘spark’.
Examples are given in \trf{03-47}.

\begin{table}
	\centering\tcp{\tsc{Flores-Lembata} final high vowels}{03-47}
	\st{6pt}\begin{tabular}{lllllll}\lsptoprule
local gloss	&	head-louse	&	stone	&	dog	&	rope	&	fire	&	pig\\
	&		&		&		&		&	*hapuy >	&	*babuy >\\
PMP	&	*kut\tbr{u}	&	*bat\tbr{u}	&	*as\tbr{u}	&	*tal\tbr{i}h	&	**ap\tbr{i} >	&	**βaβ\tbr{i}\\	\midrule
Sika	&	{\hr}ʔutu	&	{\hr}βatu	&	{\hr}ahu	&	{\hr}tali	&	{\hr\hr}api	&	{\hr\hr}βaβi\\	\midrule
Ile Mandiri	&	{\hr}kut\tbr{o}	&	{\hr}vat\tbr{o}	&	{\hr}ah\tbr{o}	&	{\hr}tal\tbr{e}ʔ	&	{\hr\hr}ap\tbr{e}	&	{\hr\hr}vav\tbr{e}\\
Adonara	&	{\hr}kut\tbr{o}	&	{\hr}wat\tbr{o}	&	{\hr}ah\tbr{o}	&	{\hr}tal\tbr{e}ʔ	&	{\hr\hr}ap\tbr{e}	&	{\hr\hr}waw\tbr{e}\\
Lebatukan 	&	{\hr}kut\tbr{o}	&	{\hr}wat\tbr{o}	&	{\hr}ah\tbr{o}	&	{\hr}tal\tbr{e}	&	{\hr\hr}ap\tbr{e}	&	\\
Lamalera	&	{\hr}kut\tbr{o}	&	{\hr}fat\tbr{o}	&	{\hr}a\tbr{o}	&	{\hr}tal\tbr{e}	&	{\hr\hr}ap\tbr{e}	&	{\hr\hr}faf\tbr{e}\\
Alorese	&	{\hr}kut\tbr{o}	&	{\hr}wat\tbr{o}	&	{\hr}ah\tbr{o}	&	{\hr}tal\tbr{e}	&	{\hr\hr}ap\tbr{e}	&	{\hr\hr}waw\tbr{e}\\
\multirow{2}{*}{Kalikasa {$\left\{\hspace{-2mm}\begin{array}{c}\vp{|}\\\vp{|}\\\end{array}\hspace{-2mm}\right.$}}
&	{\hr}kutu,	&	{\hr}vatu,	&	{\hr}au,	&	{\hr}tali,	&	{\hr\hr}api,	&	{\hr\hr}vavi,\\
&	{\hr}kut\tbr{o}-r	&	{\hr}at\tbr{o}-r	&	{\hr}a\tbr{o}-r	&	{\hr}tal\tbr{e}-r	&	{\hr\hr}ap\tbr{e}-r	&	{\hr\hr}vav\tbr{e}-r\\
Painara	&	{\hr}kut\tbr{o}-r	&	{\hr}vat\tbr{o}-r	&	{\hr}a\tbr{o}-r	&	{\hr}tali	&	{\hr\hr}ap\tbr{e}-r	&	{\hr\hr}vav\tbr{e}-r\\
Lamatuka	&	\hp{*k}ut\tbr{o}	&	{\hr}wat\tbr{o}	&	{\hr}ah\tbr{o}	&	{\hr}tal\tbr{e}	&	{\hr\hr}ap\tbr{e}	&	\\
Lewo Eleng	&	\hp{*k}ut\tbr{o}	&	{\hr}wat\tbr{o}	&	{\hr}ah\tbr{o}	&	{\hr}tal\tbr{e}	&	{\hr\hr}ap\tbr{e}	&	\\	\midrule
Kedang	&	{\hr}ʔutu	&	{\hr}waʔ	&	{\hr}au	&	\cg{(wadeʔ)}	&	{\hr\hr}api	&	{\hr\hr}wawi\\
		\lspbottomrule
	\end{tabular}
\end{table}

\begin{table}\tcp{\tsc{Flores-Lembata} *g}{03-48}
	\begin{threeparttable}\st{6pt}
	\begin{tabular}{llll}\lsptoprule
local gl.	&	itch	&	scratch	&	rough\\
PMP	&	*\tbr{g}atəl	&	*\tbr{g}aRut\su{†}	&	*\tbr{g}aRaŋ\\ \midrule
Sika	&	{\hr}\tbr{g}atar	&	{\hr}\tbr{g}aro	&	{\hr}\tbr{g}araŋ\\
Ile Mandiri 	&	{\hr}\tbr{g}atək	&	{\hr}\tbr{g}aʔu, ra\tbr{g}u	&	\\
Adonara 	&	{\hr}\tbr{g}atək	&	{\hr}\tbr{g}aru, ra\tbr{g}u	&	{\hr}\tbr{g}aʔak\\
Lamalera	&	{\hr}\tbr{g}ətər	&	{\hr}ra\tbr{g}o	&	\\
Alorese	&	{\hr}\tbr{g}ate	&	{\hr}\tbr{g}au	&	\cg{(bogal)}\\
Kalikasa	&	{\hr}\tbr{g}ətək	&	{\hr}kəra\tbr{g}u	&	\cg{(grosiŋ)}\\
Painara	&	\cg{(kənəha)}	&	{\hr}kəra\tbr{g}um	&	\\
Lamatuka	&		&	{\hr}ra\tbr{g}u	&	\\
Lewo Eleng	&		&	{\hr}ra\tbr{g}u	&	\\
\rcl Kedang	&	{\hr}\tbr{k}atel	&	{\hr}\tbr{k}aro	&	{\hr}\tbr{k}araŋ\\
		\lspbottomrule
	\end{tabular}
		\begin{tablenotes}\tnsize
			\item[†]	{Several reflexes of *gaRut have metathesis of the initial and medial consonants.}
		\end{tablenotes}
	\end{threeparttable}
\end{table}

Given the complications regarding this sound change in \tsc{Central
Lamaholot}, this sound change may only provide evidence for a node containing
\tsc{West Lamaholot} and \tsc{East Lamaholot}.
Whatever the best approach,
support for a \tsc{Lamaholot} node comes from the lexicon.
\citet{Keraf1978} shows the Lamaholot nodes are lexically closer
to one another than they are to Sika or Kedang.
The \tsc{Lamaholot} nodes share about 50{\%} shared vocabulary
with one another compared with between 20{\%}--40{\%} shared
vocabulary with Sika and Kedang.\footnote{ 
		The lowest rate of lexical similarity between members of the
		\tsc{Lamaholot} node is 44{\%} (Lewo Eleng vs.
		Lamalera and Painara vs. Pukaunu).
		The highest rate of lexical similarity for Kedang is 44{\%} with
		Ile Ape (the next highest is 40{\%} for Lewo Eleng and Lamatuka).
		The highest rate of lexical similarity for Sika is 41{\%} with
		neighbouring Pukaunu \citep[448]{Keraf1978}.}

Kedang can be placed apart from the other languages on the basis
of *g > \it{k}
and medial *d > \it{y},
though this second change may be shared with \tsc{Central Lamaholot}
as discussed above.
Examples of *g > \it{k} in Kedang are given in \trf{03-48},
and examples of medial *d > \it{y} are: *ŋuda > \it{nuyan}
`recently emerged, recently grown',
*budaq > \it{buyaʔ} `white', and *kudən > \it{uyeŋ} `pot'.

Sika can be placed apart from other languages on the basis of
*mat > \it{d} and *map > \it{b}.
Examples are given in \trf{03-49}.
Other \tsc{Flores-Lembata} languages usually have *mat > \it{t}
and *map > \it{p},
or occasionally *mat/*map > \it{m}.

\begin{table}\tcp{\tsc{Flores-Lembata} *ma-p and *ma-t}{03-49}
\begin{adjustbox}{width=\textwidth}
	\begin{threeparttable}\st{2pt}
	\begin{tabular}{lllllll}\lsptoprule
local gloss	&	ripe	&	star	&	dry	&	old	&	bitter	&	full\\
PMP	&	*\tbr{ma-t}asak	&	*ma\tbr{nt}alaq	&	*\tbr{ma-t}uquR	&	*\tbr{ma-t}uqah	&	*\tbr{ma-p}aqit	&	*\tbr{ma-p}ənuq\\ \midrule
\rcl Sika, Tana Ai	&	{\hr}\tbr{d}aha	&	{\hr}\tbr{d}ala	&	{\hr}\tbr{d}uʔur	&	{\hr}\tbr{d}uʔaŋ\su{‡}	&	{\hr}\tbr{b}aʔis	&	{\hr}\tbr{b}ənu\\
\rcl Sika, Hewa	&	{\hr}\tbr{d}ahaʔ	&	{\hr}\tbr{d}ala	&	{\hr}\tbr{d}uʔur	&	{\hr}\tbr{d}uan	&	{\hr}\tbr{b}aʔit	&	{\hr}\tbr{b}ənu\\
Ile Mandiri 	&	{\hr}\tbr{t}ahak	&	{\hr}pə\tbr{t}ala	&	{\hr}\tbr{t}uuʔ	&	{\hr}\tbr{t}\<ən\>uen	&	{\hr}\tbr{p}ait	&	{\hr}\tbr{m}ənu\\
Adonara 	&	{\hr}\tbr{t}ahak	&	{\hr}pə\tbr{t}ala, \cg{(etəp)}	&	{\hr}\tbr{t}uʔuk	&	{\hr}\tbr{t}\<ən\>uʔe	&	{\hr}\tbr{p}ait	&	{\hr}\tbr{p}əno\\
Alorese	&	{\hr}\tbr{t}ahak	&	{\hr}\tbr{t}əmala\su{†}	&	\cg{(marak)}	&	{\hr}\tbr{t}ua	&	{\hr}\tbr{p}ai	&	{\hr}\tbr{p}ənoŋ\\
Kalikasa	&	{\hr}\tbr{t}asak	&	\cg{(tona)}	&	\cg{(maʤak)}	&	{\hr}\tbr{t}uana	&	{\hr}\tbr{p}‹ən›ait	&	{\hr}\tbr{m}ənuk\\
Lamatuka	&		&	{\hr}malã	&	\cg{(marʔã)}	&	\cg{(magũ)}	&		&	\\
Kedang	&	{\hr}\tbr{t}aʔen	&	{\hr}\tbr{t}ala, male \tbr{t}ala	&	\cg{(mayaʔ)}	&	{\hr}\tbr{t}ua lahar	&	{\hr}\tbr{p}æiʔ	&	{\hr}\tbr{p}ænu\\
		\lspbottomrule
	\end{tabular}
		\begin{tablenotes}\tnsize
			\item[†]	{Alorese \it{təmala} has consonant metathesis from intermediate **mətala.
								Lamatuka \it{malã} and Kedang \it{male} are probably from **təmala
								with loss of the antepenultimate syllable.}
			\item[‡]	{Sika \it{kobar duʔaŋ} is ‘half ripened coconut’.}
		\end{tablenotes}
	\end{threeparttable}
\end{adjustbox}
\end{table}

Regarding final vowel-glide sequences,
Sika has *-aw > \it{o}
and *-ay > \it{e}.
Kedang has *-aw > \it{o} (when retained)
and *ay > \it{e, a}.
\tsc{Central Lamaholot} usually has *-aw > \it{a,
av} and *-ay > \it{a,
ey, aʤ} (variety dependent).
\tsc{West Lamaholot} has *-aw > \it{a},
and a variety of reflexes for *-ay including: \it{i, e,} and \it{a}.
Examples are given in \trf{03-50}.
While the reflexes are diverse,
they show that the full vowel-glide sequence
needs to be reconstructed to Proto-Flores-Lembata.

\begin{table}\tcp{\tsc{Flores-Lembata} *-aw and *-ay}{03-50}
\begin{adjustbox}{width=\textwidth}
	\begin{threeparttable}\st{2pt}
	\begin{tabular}{llll|lll}\lsptoprule
local gloss	&	walk, go	&	steal	&	sun	&	die	&	rattan	&	liver\\
PMP	&	*pan\tbr{aw}	&	*ma-tak\tbr{aw}	&	*qaləj\tbr{aw}	&	*mat\tbr{ay}	&	*qu\tbr{ay}	&	*qat\tbr{ay}\\	\midrule
Sika, Tana Ai	&	{\hr}pan\tbr{o}	&	{\hr}naʔ\tbr{o}	&	{\hqa}lər\tbr{o}	&	{\hr}mat\tbr{e}	&	\cg{(gai)}	&	{\hr}ʔβat\tbr{e}-ŋ\\
Sika, Hewa	&	{\hr}pan\tbr{o}	&	\cg{(toʔi)}	&	{\hqa}lər\tbr{o}	&	{\hr}mat\tbr{e}	&	{\hr}u\tbr{e}	&	{\hr}vat\tbr{e}-n\\
Pukaunu	&	{\hr}pan\tbr{a}	&		&	{\hqa}lər\tbr{a}	&	{\hr}mat\tbr{ĩ}	&		&	{\hr}at\tbr{e}\\
Lewoingu\su{†}	&	{\hr}pan\tbr{a}	&	{\hr}təmak\tbr{a}	&	{\hqa}lər\tbr{a}	&	{\hr}mat\tbr{a}	&	{\hr}uw\tbr{ay}əŋ	&	{\hr}at\tbr{e}-ŋ\\
Lewolema\su{†} 	&	{\hr}pan\tbr{a}	&	{\hr}tak\tbr{a}ʔ, təmak\tbr{a}ʔ	
					&	{\hqa}rər\tbr{a}	&	{\hr}mat\tbr{a}	&	{\hr}u\tbr{a}, u\tbr{ar}	&	{\hr}at\tbr{e}\\
Adonara 	&	{\hr}pan\tbr{a}	&	\cg{(lãã)}	&	{\hqa}rər\tbr{a}	&	{\hr}mat\tbr{a}, mat\tbr{e}	&	{\hr}u\tbr{a}	&	{\hr}at\tbr{e}\\
Alorese	&	{\hr}pan\tbr{a}	&	{\hr}təmak\tbr{a}, mak\tbr{a}	
					&	{\hqa}lər\tbr{a}	&	{\hr}mat\tbr{e}	&	{\hr}uw\tbr{e}	&	{\hr}at\tbr{e}ŋ\\
Kalikasa	&	{\hr}pan\tbr{a}, 
		pan\tbr{av}	&	{\hr}tak\tbr{av}	&	\hp{qa}\cg{(luvak)}	&	{\hr}mat\tbr{e}k, mat\tbr{aʤ}	&	{\hr}u\tbr{a}, u\tbr{aʤ}	&	\cg{(oren)}\\
Painara	&	{\hr}pan\tbr{ə}	&	{\hr}tak\tbr{av}	&	{\hqa}ləʤ\tbr{aw}	&	{\hr}mat\tbr{ey}	&	{\hr}u\tbr{ar}	&	\cg{(orən)}\\
Labalekan	&	{\hr}pan\tbr{a}	&		&	{\hqa}ləʤ\tbr{af}	&	{\hr}mat\tbr{aʤ}	&		&	\cg{(oraha)}\\
Lamatuka	&	{\hr}pan\tbr{a}	&		&	\hp{qa}\cg{(luwa)}	&	{\hr}mat\tbr{a}	&		&	\cg{(huʔa)}\\
Lewo Eleng	&	{\hr}pã\tbr{ã}	&		&	{\hqa}lər\tbr{a}	&	{\hr}mat\tbr{a}	&		&	\cg{(ũhã)}\\
Kedang	&	{\hr}pan	&	{\hr}maʔ\tbr{o}	&	{\hqa}loy\tbr{o}	&	{\hr}mat\tbr{e}	&	{\hr}u\tbr{a}	&	\cg{(oneʔ)}\\
		\lspbottomrule
	\end{tabular}
		\begin{tablenotes}\tnsize
			\item[†]	{Lewoingu and Lewolema are both dialects/varieties of Ile Mandiri Lamaholot [slp].}
		\end{tablenotes}
	\end{threeparttable}
\end{adjustbox}
\end{table}
